\subsection{Berechnung des Vorwiderstands einer LED}
% Alt: Aufgabe 10

\Aufgabe{
    \label{Aufg10}
    Wie bei allen Dioden muss auch bei Leuchtdioden (LEDs) der Strom begrenzt werden. Dazu wird ein Vorwiderstand in Reihe geschaltet, wie in der unten gezeigten Schaltung. \newline
    Die eingesetzte LED besitzt eine Durchlassspannung von $U_\mathrm{D} = 2,2\,\mathrm{V}$. Die ideale Helligkeit wird bei einem Strom von $I_\mathrm{D} = 30\,\mathrm{mA}$ erreicht. Die Eingangsspannung beträgt $U_\mathrm{0} = 5\,\mathrm{V}$.

    \begin{figure}[H]
        \centering
        \input{Tikz/src/FigLEDVorwiderstand.tex}\alttext{Zu sehen ist eine elektrische Schaltung mit drei Bauteilen zur Bestimmung des Vorwiderstands einer Leuchtdiode. Auf der linken Seite befindet sich die Spannungsquelle \(U_0\), deren Spannungspfeil nach unten zeigt. In Reihe zur Spannungsquelle befindet sich  ein Vorwiderstand \(R_\mathrm{V}\) und eine Leuchtdiode (LED). Neben der Diode befindet sich zusätzlich ein Spannungspfeil in Stromrichtung mit der Beschriftung \(U_\mathrm{D}\). Zusätzlich wird der Stromfluss in Uhrzeigerrichtung durch einen roten Pfeil mit der Beschriftung \(I_\mathrm{D}\) symbolisiert.}
        \label{fig:FigLEDVorwiderstand}
    \end{figure}

    Welchen der folgenden Widerstandswerte würden Sie für $R$ verwenden, um die LED möglichst nahe an ihrem optimalen Betriebspunkt (d.h. mit optimaler Helligkeit) zu betreiben?\\
    $15\,\Omega$, $33\,\Omega$, $68\,\Omega$, $150\,\Omega$, $180\,\Omega$
}


\Loesung{
    Berechnung der Spannung am Vorwiderstand:
    \begin{equation*}
        U_\mathrm{V} = U_\mathrm{0} - U_\mathrm{D} = 5\,\mathrm{V} - 2,2\,\mathrm{V} = 2,8\,\mathrm{V}
    \end{equation*}

    Berechnung des optimalen Vorwiderstands für \( I_\mathrm{D} = \mathrm{30\,mA} \):
    \begin{equation*}
        R_\mathrm{V} = \frac{U_\mathrm{V}}{I_\mathrm{D}} = \frac{2,8\,\mathrm{V}}{30\,\mathrm{mA}} = 93,3\,\Omega
    \end{equation*}

    Berechnung des Stroms bei Auswahl eines konkreten Widerstandswertes:
    \begin{equation*}
        I_\mathrm{D} = \frac{U_\mathrm{V}}{R} = \frac{2,8\,\mathrm{V}}{150\,\Omega} = 18,7\,\mathrm{mA}
    \end{equation*}

    Berechnung bei zwei Widerständen in Parallelschaltung mit zusätzlichem Serienwiderstand:
    \begin{gather*}
        R_\mathrm{ges} = \frac{150\,\Omega \cdot 180\,\Omega}{150\,\Omega + 180\,\Omega} + 15\,\Omega = 96,8\,\Omega \\
        I_\mathrm{D} = \frac{2,8\,\mathrm{V}}{96,8\,\Omega} =28,9\,\mathrm{mA}
    \end{gather*}

    Alternative Reihenschaltung:
    \begin{equation*}
        R_\mathrm{ges} = 4 \cdot 15\,\Omega + 33\,\Omega = 93\,\Omega
    \end{equation*}

    \textbf{Fazit:}
    Ein kombinierter Widerstand aus $4 \cdot 15\,\Omega + 33\,\Omega = 93\,\Omega$ oder eine Parallelschaltung mit $150\,\Omega \,\|\, 180\,\Omega + 15\,\Omega = 96,8\,\Omega$ führt zu einem Strom nahe dem optimalen Wert von  $30\,\mathrm{mA}$. \\
    \textbf{Empfohlene Wahl:} Kombination mit Gesamtwiderstand nahe $93\,\Omega$.

    Siehe: Abschnitt \ref{sec:SpezielleDioden}
}